%% ============================================================
%% COMM514 MSc Report -- Section 4: System Implementation
%% University of Exeter
%% ============================================================
%% Embed in your main document with %% ============================================================
%% COMM514 MSc Report -- Section 4: System Implementation
%% University of Exeter
%% ============================================================
%% Embed in your main document with %% ============================================================
%% COMM514 MSc Report -- Section 4: System Implementation
%% University of Exeter
%% ============================================================
%% Embed in your main document with %% ============================================================
%% COMM514 MSc Report -- Section 4: System Implementation
%% University of Exeter
%% ============================================================
%% Embed in your main document with \input{implementation}
%% ============================================================

\section{System Implementation}

\subsection{Overview}

The experiment requires two functionally complete code generation systems. The first, the baseline, is the control condition: a minimal single-agent pipeline in which a participant submits a coding task and receives GPT-4o's response directly, without any intermediate processing. The second, the multi-agent system, is the experimental condition: a five-agent pipeline orchestrated through LangGraph in which each agent performs a discrete, specialised function before the next begins. Both systems use GPT-4o at a fixed temperature of 0.2, are served by the same FastAPI backend, and present to participants through the same Next.js frontend application. The sole variable between conditions is architectural structure. Figure~\ref{fig:architecture} shows the complete deployed system architecture.

\subsection{Baseline Agent}

The baseline agent is implemented in \textit{baseline/agent.py} using the OpenAI Python SDK directly, without any orchestration framework. When a participant submits a task, the \textit{run\_baseline(task, client)} function constructs a two-message prompt containing a brief system message and the task as the user turn, and submits a single completion request to GPT-4o with \texttt{max\_tokens=2048}. The system prompt is deliberately minimal:

\begin{quote}
\textit{You are a software engineer assistant. When given a coding task, produce working, readable code. Do not add commentary beyond what is needed to understand the code. Do not omit error handling.}
\end{quote}

The prompt contains no security instructions, no structured review steps, and no defensive coding guidance. This is intentional: any security guidance in the baseline prompt would reduce the architectural contrast between conditions and bias the comparison. The model response is returned to the participant without further processing.

Session data is written to \textit{logs/baseline\_sessions.jsonl} by \textit{log\_session()}, which appends one JSON object per session containing the fields: \texttt{timestamp} (ISO 8601 UTC), \texttt{condition} (fixed value \texttt{"baseline"}), \texttt{participant\_id}, \texttt{model}, \texttt{temperature}, \texttt{task}, \texttt{task\_id}, \texttt{task\_order}, \texttt{response}, \texttt{duration\_seconds}, and an optional \texttt{participant\_notes} field. The \textit{run\_baseline} function and \textit{log\_session} function are importable module-level callables. The interactive command-line interface for development testing is confined to the \texttt{if \_\_name\_\_ == "\_\_main\_\_"} block and is never imported by the application.

The OpenAI client, model identifier, and temperature are imported from \textit{config.py}, which reads the API key from the environment via \texttt{python-dotenv}. No API credentials are instantiated inside agent files or committed to version control.

\subsection{Multi-Agent System}

The multi-agent system is built on LangGraph's \textit{StateGraph} class, which represents the pipeline as a directed acyclic graph in which each node is an asynchronous function that reads from a shared state object and returns a partial state update. The shared state schema is defined in \textit{multiagent/state.py} as the \textit{AgentState} \textit{TypedDict}, which contains all fields required by all five agents, organised by pipeline stage. Every agent output field is typed as \textit{Optional} because the state is populated progressively and earlier agents cannot know what later agents will produce.

The graph is assembled in \textit{multiagent/graph.py}. Six nodes are registered: \textit{planner}, \textit{threat\_modeller}, \textit{code\_generator}, \textit{code\_reviewer}, \textit{verifier}, and \textit{finalise}. Edges connect them in a linear sequence. The graph is compiled with a \textit{MemorySaver} checkpointer, which persists state across asynchronous resume calls, and with \texttt{interrupt\_after=["planner", "threat\_modeller", "code\_generator", "code\_reviewer"]}. This configuration pauses execution after each of the first four nodes, allowing the frontend to display that node's output to the participant and collect their decision before the pipeline continues.

Three public functions expose the graph to the API layer. \textit{start\_session()} initialises an \textit{AgentState}, generates a unique thread identifier from the participant ID and a UTC timestamp, runs the graph from the entry point through the first interrupt, and returns the thread configuration. \textit{resume\_session()} applies a participant decision as a state update and resumes execution from the next node. \textit{get\_current\_state()} retrieves the latest state snapshot for a given thread, which the frontend polls to read agent output between interrupts.

When the Verifier completes, the graph runs the \textit{\_finalise} node automatically without an additional interrupt. This node computes the total session duration, calls \textit{to\_log\_entry()} to assemble the JSONL record, and appends it to \textit{logs/multiagent\_sessions.jsonl}. The top-level fields of the multi-agent log record (\texttt{timestamp}, \texttt{condition}, \texttt{participant\_id}, \texttt{model}, \texttt{temperature}, \texttt{task}, \texttt{task\_id}, \texttt{task\_order}, \texttt{response}, \texttt{duration\_seconds}) are identical to the baseline schema. An additional \texttt{multi\_agent\_detail} field contains the per-agent outputs, the participant's decision at each interrupt point, and the full HITL coding metrics collected during the code-writing stage.

\subsection{Agent Implementations}

\subsubsection*{Planner}

The Planner receives the participant's coding task from state and produces a structured implementation plan. Its system prompt instructs GPT-4o to decompose the task into three to six concrete, ordered implementation steps, a single-sentence scope statement describing what the code should and should not do, and a list of three to five security properties the final code must satisfy. The prompt explicitly distinguishes concrete security requirements (such as "Hash passwords with bcrypt; do not compare plaintext") from generic advice. The API call uses \texttt{response\_format=\{"type": "json\_object"\}} to guarantee that the response is valid JSON, eliminating the need for error-prone string parsing. The parsed output is stored in \texttt{state["plan"]} as a \textit{PlannerOutput} dictionary. The Planner uses no external tools.

\subsubsection*{Threat Modeller}

The Threat Modeller receives the task and the Planner output and produces a structured threat model before any code is written. It executes three operations in sequence. First, it calls \textit{retrieve(task)} from \textit{rag/retriever.py}, which runs the full four-stage Advanced RAG pipeline described in Section 4.5, returning the three most relevant CWE entries from the local corpus. Second, it queries the NIST NVD REST API v2 via the \textit{search\_nvd} MCP tool, using the short name of the top-ranked CWE entry as the search keyword and requesting three CVEs; the call is wrapped in \texttt{asyncio.wait\_for} with a 25-second timeout, so a slow or unavailable NVD service causes graceful degradation to CWE-only context rather than a pipeline failure. Third, it assembles the task, plan scope, Planner security requirements, formatted CWE context, and formatted CVE context into a single user message and calls GPT-4o to produce a JSON list of threats. Each threat is a \textit{ThreatEntry} dictionary with fields \texttt{cwe\_id}, \texttt{name}, \texttt{severity}, \texttt{description} (task-specific), and \texttt{mitigation} (a direct, implementable instruction for the Code Generator). The raw RAG retrieval results are stored in \texttt{state["rag\_context"]} for logging transparency.

\subsubsection*{Code Generator}

The Code Generator node implements the most distinctive architectural feature of the multi-agent system: a human-in-the-loop coding mentor. In this design, the participant writes all code themselves; no AI-generated code is produced or presented for acceptance. The LangGraph node function \textit{run\_code\_generator()} sets \texttt{state["current\_stage"]} to \texttt{"coding\_in\_progress"} and returns immediately, signalling to the graph that the pipeline is waiting for the human participant rather than completing an autonomous step.

During the coding stage, the frontend provides the participant with four asynchronous hint functions. \textit{get\_step\_hint(state, step\_index, level)} returns a hint for a specific plan step at one of four escalating levels: \textit{direction} (plain English explaining what the step requires and why); \textit{pseudocode} (comment-only skeleton with no executable code); \textit{partial} (structural implementation with \texttt{\# TODO} markers for security-critical logic); and \textit{full} (complete implementation for boilerplate steps). Each level is defined by its own system prompt with strict instructions about what may and may not be revealed. \textit{get\_next\_hint(state, code\_so\_far)} reads the participant's current code and returns a focused suggestion identifying which plan step appears to need attention next, without rewriting any code. \textit{get\_security\_hint(state, code\_so\_far)} scans the current code against the threat model and flags the single most important security issue or missing mitigation.

The maximum hint level available to the participant for each plan step is not uniform. Prior to the coding stage, a step classifier (\textit{multiagent/step\_classifier.py}) uses a separate GPT-4o call at temperature 0.0 to read the plan steps and the threat model mitigations together, assigning each step a cap of \textit{direction}, \textit{pseudocode}, \textit{partial}, or \textit{full} based on whether the step directly implements a named security mitigation, is security-adjacent, is general logic, or is pure boilerplate respectively. Steps classified as \textit{direction} receive only plain-English guidance regardless of what the participant requests; the corresponding code is never revealed. The per-step caps are stored in \texttt{state["step\_hint\_caps"]} and enforced by the API route handler before any hint function is called.

When the participant is ready to submit their code, the frontend calls \textit{finalize\_code(state, user\_code, annotations, confidence\_rating, hints\_requested, time\_in\_coding\_seconds)}, which stores the participant's code in \texttt{state["generated\_code"]} and packages the full HITL metrics into a \textit{HitlCodingMetrics} dictionary: the deepest hint level reached across all steps (as a numeric depth score of one to four), the ordered list of all hint requests with timestamps, total time in seconds spent at the coding stage, a self-assessed confidence rating from one to five, and the annotations the participant provided at the submission gate. The annotation gate requires the participant to state in free text what their code does and to list which threats they believe they have addressed, before the Code Reviewer is permitted to run.

\subsubsection*{Code Reviewer}

The Code Reviewer receives the participant's code and the threat model from state and conducts an independent security review in two sequential stages. First, it calls the \textit{run\_bandit} MCP tool, which executes Bandit static analysis on the code and returns structured JSON findings. Second, it calls GPT-4o with the code, the formatted Bandit findings, and the threat model mitigations as combined context, instructing the model to include all Bandit findings in the output (tagged \texttt{"source": "bandit"}), add any further issues Bandit missed (tagged \texttt{"source": "llm"}), and verify whether each threat model mitigation was correctly implemented. Each finding in the output is a \textit{ReviewFinding} dictionary with fields \texttt{cwe\_id}, \texttt{severity}, \texttt{description} (referencing specific lines), \texttt{suggested\_fix} (a concrete, implementable fix), \texttt{line\_number}, and \texttt{source}. The raw Bandit findings are stored in \texttt{state["bandit\_findings\_review"]}, which is a separate state field from the Verifier's Bandit run. These two fields are never merged; the Verifier is a genuine independent check, not a repeat.

\subsubsection*{Verifier}

The Verifier performs the final independent validation of the participant's code against the original threat model. It executes four operations. First, it calls the \textit{run\_bandit} MCP tool independently, storing results in \texttt{state["bandit\_findings\_verify"]}. Second, it calls the \textit{execute\_code} MCP tool, which runs the code in a sandboxed subprocess and returns a pass or fail result with exit code, stdout, and stderr. Third, it calls \textit{retrieve(task)} on the CWE corpus independently, with no access to the Threat Modeller's earlier retrieval results, arriving at its own security context. Fourth, it calls GPT-4o with the code, the threat model mitigations, the independent Bandit findings, the independent CWE context, and the sandbox execution result, instructing the model to produce a per-threat verdict (\textit{passed}: true only if the mitigation is fully and correctly implemented, \textit{notes}: a concrete observation citing the specific line or pattern observed). The overall pass verdict is true only when all per-threat verdicts pass and the sandbox execution returns exit code zero.

\subsection{RAG Pipeline}

The Retrieval-Augmented Generation pipeline provides both the Threat Modeller and the Verifier with structured, CWE-specific security context drawn from the MITRE CWE Top 25 Most Dangerous Software Weaknesses corpus. The corpus is stored as a structured JSON file at \textit{rag/cwe\_corpus/cwe\_top25.json}.

Ingestion is handled by \textit{rag/ingest.py}. Each of the 25 CWE entries is split into three focused chunks by \textit{chunk\_entry()}: a main description chunk (containing the CWE identifier, rank, domain, severity, full description, and extended description); a mitigations chunk (containing the prevention strategies for this weakness); and an examples chunk (containing vulnerable code scenarios). This hierarchical splitting ensures that a query about how to prevent a weakness retrieves the mitigations chunk specifically, rather than a general entry in which mitigations constitute only a fraction of the text. The 75 chunks in total are embedded using the OpenAI \texttt{text-embedding-3-small} model and stored in a local ChromaDB persistent collection configured with cosine distance. Each chunk carries metadata fields for \texttt{cwe\_id}, \texttt{rank}, \texttt{severity}, \texttt{domain}, and \texttt{chunk\_type}, which the retriever uses for filtering.

Retrieval is handled by \textit{rag/retriever.py}, which implements the four-stage pipeline. In stage one, \textit{rewrite\_query(task)} calls GPT-4o to expand the coding task into security-relevant search terminology and to identify which of seven available domain labels (\textit{injection}, \textit{memory}, \textit{access\_control}, \textit{input\_validation}, \textit{information\_disclosure}, \textit{network}, \textit{resource\_management}) are relevant to the task. In stage two, ChromaDB's \texttt{where} filter restricts the vector search to chunks whose \texttt{domain} field matches one of the identified domains, reducing the candidate set before any embedding comparison is performed. In stage three, ChromaDB returns the ten closest chunks by cosine similarity to the expanded search query. In stage four, \textit{rerank(task, candidates)} calls GPT-4o to select the three most relevant candidates from the ten returned by the similarity search, reasoning about task fit rather than vector distance alone. The final output is deduplicated by CWE identifier so that the three returned chunks cover distinct weaknesses. The formatted context string returned by \textit{format\_for\_prompt()} is injected directly into each consuming agent's user message.

\subsection{MCP Servers}

The three Model Context Protocol servers are implemented using the \texttt{mcp.server.Server} class from the \texttt{mcp} library, each running over a \texttt{stdio\_server()} transport. They are connected to the agent layer via \textit{mcp\_client.py}, which configures a \textit{MultiServerMCPClient} from \texttt{langchain-mcp-adapters} to manage all three servers as subprocess connections. Each server is launched using the current Python interpreter (\texttt{sys.executable}) so it inherits the project virtual environment. A single \textit{call\_tool(client, tool\_name, arguments)} helper function provides a uniform interface for all agent tool calls, resolving the tool by name across the combined flat tool list and parsing the JSON response into a Python dictionary.

\textit{mcp\_servers/bandit\_server.py} exposes a single tool, \textit{run\_bandit(code)}. The implementation writes the code string to a temporary file using \texttt{tempfile.NamedTemporaryFile}, executes \texttt{bandit -f json -q} on that file as a subprocess with a 30-second timeout, parses the JSON output, normalises each finding into a dictionary with \texttt{test\_id}, \texttt{test\_name}, \texttt{severity}, \texttt{confidence}, \texttt{description}, \texttt{line\_number}, \texttt{cwe\_id}, and \texttt{code\_snippet} fields, and deletes the temporary file in a \texttt{finally} block regardless of whether an error occurred. Bandit exit code zero indicates no issues found; exit code one indicates issues found normally; exit code two indicates that Bandit itself failed to run and is treated as an error.

\textit{mcp\_servers/nist\_nvd\_server.py} exposes a single tool, \textit{search\_nvd(keyword, max\_results)}. The implementation uses \texttt{httpx.AsyncClient} with a ten-second request timeout to query the NVD REST API v2 endpoint at \texttt{services.nvd.nist.gov/rest/json/cves/2.0}, normalises the nested NVD response into a flat list of CVE dictionaries with English description, CVSS v3 severity and score, publication date, associated CWE identifiers, and NVD URL, and returns up to ten results. HTTP 403 and 429 rate-limit responses are handled gracefully with informative error messages. An optional \texttt{NVD\_API\_KEY} environment variable is supported to raise the rolling rate limit from five to fifty requests per thirty seconds.

\textit{mcp\_servers/sandbox\_server.py} exposes a single tool, \textit{execute\_code(code)}. Before spawning any subprocess, the implementation validates the code string with \texttt{ast.parse}, returning a clean syntax-error response immediately if parsing fails. Valid code is written to an isolated temporary directory using \texttt{tempfile.mkdtemp}, which confines the execution working directory to a path outside the project folder so relative file path references cannot reach project files. The code is executed using \texttt{asyncio.create\_subprocess\_exec} with \texttt{stdin=DEVNULL}, integrating the subprocess into the MCP server's event loop and preventing the child process from inheriting the parent's MCP stdio pipe. Execution is bounded by a ten-second timeout via \texttt{asyncio.wait\_for}; if the timeout expires, the process is killed and the \texttt{timed\_out} field in the response is set to true. Stdout and stderr are each capped at 4,096 bytes to prevent memory exhaustion from runaway output. The temporary directory is removed in a \texttt{finally} block.

\subsection{Frontend and Session Management}

The participant-facing interface is a Next.js 15 application using React 19 and TypeScript. Both study views (\textit{frontend/app/study/baseline/page.tsx} and \textit{frontend/app/study/multiagent/page.tsx}) use the Monaco Editor component for code display and editing, loaded asynchronously to avoid blocking the initial page render.

The baseline view manages a five-phase flow: loading (fetching a shuffled task list from the API and reading the participant identifier from session storage), ready (displaying the task and an editor where the participant reads the AI-generated code and may annotate it), generating (awaiting the API response), reviewing (displaying the result with an annotation text area for the participant's notes), and complete (logging the session and advancing to the next task). Task order is randomised server-side on first fetch and the shuffled sequence is stored in session storage so it persists across page refreshes within a session.

The multi-agent view manages a more complex stage-detection flow derived directly from the pipeline state object returned by the API. A \textit{detectStage(state)} function maps the combination of state fields present (plan approved or not, threats approved or not, current\_stage value, review findings present or not) to one of the named UI stages: \textit{plan\_review}, \textit{threat\_review}, \textit{coding}, \textit{code\_review}, \textit{task\_complete}, or \textit{agent\_error}. During the coding stage, the frontend exposes the step plan in a sidebar, calls \textit{fetchStepCaps} to retrieve the per-step hint level caps from the step classifier, and presents hint request controls scoped to each plan step. The \textit{requestStepHint}, \textit{requestNextHint}, and \textit{requestSecurityHint} API calls each invoke the corresponding Code Generator functions. When the participant is ready to submit, the frontend collects their annotation gate responses, confidence rating, and full hint history and calls \textit{finalizeCode}, which resumes the graph and triggers the Code Reviewer. After the Code Reviewer output is displayed, the participant may call \textit{reviseCode} to update their submission before approving the review and releasing the Verifier.

Participant identifiers are assigned at session start, read from URL parameters set by the study coordinator, and stored in browser session storage for the duration of the session. All JSONL log records are written server-side; no sensitive session data is stored in the browser beyond the identifier and task sequence needed to reconstruct the current session state.

%% ============================================================

\section{System Implementation}

\subsection{Overview}

The experiment requires two functionally complete code generation systems. The first, the baseline, is the control condition: a minimal single-agent pipeline in which a participant submits a coding task and receives GPT-4o's response directly, without any intermediate processing. The second, the multi-agent system, is the experimental condition: a five-agent pipeline orchestrated through LangGraph in which each agent performs a discrete, specialised function before the next begins. Both systems use GPT-4o at a fixed temperature of 0.2, are served by the same FastAPI backend, and present to participants through the same Next.js frontend application. The sole variable between conditions is architectural structure. Figure~\ref{fig:architecture} shows the complete deployed system architecture.

\subsection{Baseline Agent}

The baseline agent is implemented in \textit{baseline/agent.py} using the OpenAI Python SDK directly, without any orchestration framework. When a participant submits a task, the \textit{run\_baseline(task, client)} function constructs a two-message prompt containing a brief system message and the task as the user turn, and submits a single completion request to GPT-4o with \texttt{max\_tokens=2048}. The system prompt is deliberately minimal:

\begin{quote}
\textit{You are a software engineer assistant. When given a coding task, produce working, readable code. Do not add commentary beyond what is needed to understand the code. Do not omit error handling.}
\end{quote}

The prompt contains no security instructions, no structured review steps, and no defensive coding guidance. This is intentional: any security guidance in the baseline prompt would reduce the architectural contrast between conditions and bias the comparison. The model response is returned to the participant without further processing.

Session data is written to \textit{logs/baseline\_sessions.jsonl} by \textit{log\_session()}, which appends one JSON object per session containing the fields: \texttt{timestamp} (ISO 8601 UTC), \texttt{condition} (fixed value \texttt{"baseline"}), \texttt{participant\_id}, \texttt{model}, \texttt{temperature}, \texttt{task}, \texttt{task\_id}, \texttt{task\_order}, \texttt{response}, \texttt{duration\_seconds}, and an optional \texttt{participant\_notes} field. The \textit{run\_baseline} function and \textit{log\_session} function are importable module-level callables. The interactive command-line interface for development testing is confined to the \texttt{if \_\_name\_\_ == "\_\_main\_\_"} block and is never imported by the application.

The OpenAI client, model identifier, and temperature are imported from \textit{config.py}, which reads the API key from the environment via \texttt{python-dotenv}. No API credentials are instantiated inside agent files or committed to version control.

\subsection{Multi-Agent System}

The multi-agent system is built on LangGraph's \textit{StateGraph} class, which represents the pipeline as a directed acyclic graph in which each node is an asynchronous function that reads from a shared state object and returns a partial state update. The shared state schema is defined in \textit{multiagent/state.py} as the \textit{AgentState} \textit{TypedDict}, which contains all fields required by all five agents, organised by pipeline stage. Every agent output field is typed as \textit{Optional} because the state is populated progressively and earlier agents cannot know what later agents will produce.

The graph is assembled in \textit{multiagent/graph.py}. Six nodes are registered: \textit{planner}, \textit{threat\_modeller}, \textit{code\_generator}, \textit{code\_reviewer}, \textit{verifier}, and \textit{finalise}. Edges connect them in a linear sequence. The graph is compiled with a \textit{MemorySaver} checkpointer, which persists state across asynchronous resume calls, and with \texttt{interrupt\_after=["planner", "threat\_modeller", "code\_generator", "code\_reviewer"]}. This configuration pauses execution after each of the first four nodes, allowing the frontend to display that node's output to the participant and collect their decision before the pipeline continues.

Three public functions expose the graph to the API layer. \textit{start\_session()} initialises an \textit{AgentState}, generates a unique thread identifier from the participant ID and a UTC timestamp, runs the graph from the entry point through the first interrupt, and returns the thread configuration. \textit{resume\_session()} applies a participant decision as a state update and resumes execution from the next node. \textit{get\_current\_state()} retrieves the latest state snapshot for a given thread, which the frontend polls to read agent output between interrupts.

When the Verifier completes, the graph runs the \textit{\_finalise} node automatically without an additional interrupt. This node computes the total session duration, calls \textit{to\_log\_entry()} to assemble the JSONL record, and appends it to \textit{logs/multiagent\_sessions.jsonl}. The top-level fields of the multi-agent log record (\texttt{timestamp}, \texttt{condition}, \texttt{participant\_id}, \texttt{model}, \texttt{temperature}, \texttt{task}, \texttt{task\_id}, \texttt{task\_order}, \texttt{response}, \texttt{duration\_seconds}) are identical to the baseline schema. An additional \texttt{multi\_agent\_detail} field contains the per-agent outputs, the participant's decision at each interrupt point, and the full HITL coding metrics collected during the code-writing stage.

\subsection{Agent Implementations}

\subsubsection*{Planner}

The Planner receives the participant's coding task from state and produces a structured implementation plan. Its system prompt instructs GPT-4o to decompose the task into three to six concrete, ordered implementation steps, a single-sentence scope statement describing what the code should and should not do, and a list of three to five security properties the final code must satisfy. The prompt explicitly distinguishes concrete security requirements (such as "Hash passwords with bcrypt; do not compare plaintext") from generic advice. The API call uses \texttt{response\_format=\{"type": "json\_object"\}} to guarantee that the response is valid JSON, eliminating the need for error-prone string parsing. The parsed output is stored in \texttt{state["plan"]} as a \textit{PlannerOutput} dictionary. The Planner uses no external tools.

\subsubsection*{Threat Modeller}

The Threat Modeller receives the task and the Planner output and produces a structured threat model before any code is written. It executes three operations in sequence. First, it calls \textit{retrieve(task)} from \textit{rag/retriever.py}, which runs the full four-stage Advanced RAG pipeline described in Section 4.5, returning the three most relevant CWE entries from the local corpus. Second, it queries the NIST NVD REST API v2 via the \textit{search\_nvd} MCP tool, using the short name of the top-ranked CWE entry as the search keyword and requesting three CVEs; the call is wrapped in \texttt{asyncio.wait\_for} with a 25-second timeout, so a slow or unavailable NVD service causes graceful degradation to CWE-only context rather than a pipeline failure. Third, it assembles the task, plan scope, Planner security requirements, formatted CWE context, and formatted CVE context into a single user message and calls GPT-4o to produce a JSON list of threats. Each threat is a \textit{ThreatEntry} dictionary with fields \texttt{cwe\_id}, \texttt{name}, \texttt{severity}, \texttt{description} (task-specific), and \texttt{mitigation} (a direct, implementable instruction for the Code Generator). The raw RAG retrieval results are stored in \texttt{state["rag\_context"]} for logging transparency.

\subsubsection*{Code Generator}

The Code Generator node implements the most distinctive architectural feature of the multi-agent system: a human-in-the-loop coding mentor. In this design, the participant writes all code themselves; no AI-generated code is produced or presented for acceptance. The LangGraph node function \textit{run\_code\_generator()} sets \texttt{state["current\_stage"]} to \texttt{"coding\_in\_progress"} and returns immediately, signalling to the graph that the pipeline is waiting for the human participant rather than completing an autonomous step.

During the coding stage, the frontend provides the participant with four asynchronous hint functions. \textit{get\_step\_hint(state, step\_index, level)} returns a hint for a specific plan step at one of four escalating levels: \textit{direction} (plain English explaining what the step requires and why); \textit{pseudocode} (comment-only skeleton with no executable code); \textit{partial} (structural implementation with \texttt{\# TODO} markers for security-critical logic); and \textit{full} (complete implementation for boilerplate steps). Each level is defined by its own system prompt with strict instructions about what may and may not be revealed. \textit{get\_next\_hint(state, code\_so\_far)} reads the participant's current code and returns a focused suggestion identifying which plan step appears to need attention next, without rewriting any code. \textit{get\_security\_hint(state, code\_so\_far)} scans the current code against the threat model and flags the single most important security issue or missing mitigation.

The maximum hint level available to the participant for each plan step is not uniform. Prior to the coding stage, a step classifier (\textit{multiagent/step\_classifier.py}) uses a separate GPT-4o call at temperature 0.0 to read the plan steps and the threat model mitigations together, assigning each step a cap of \textit{direction}, \textit{pseudocode}, \textit{partial}, or \textit{full} based on whether the step directly implements a named security mitigation, is security-adjacent, is general logic, or is pure boilerplate respectively. Steps classified as \textit{direction} receive only plain-English guidance regardless of what the participant requests; the corresponding code is never revealed. The per-step caps are stored in \texttt{state["step\_hint\_caps"]} and enforced by the API route handler before any hint function is called.

When the participant is ready to submit their code, the frontend calls \textit{finalize\_code(state, user\_code, annotations, confidence\_rating, hints\_requested, time\_in\_coding\_seconds)}, which stores the participant's code in \texttt{state["generated\_code"]} and packages the full HITL metrics into a \textit{HitlCodingMetrics} dictionary: the deepest hint level reached across all steps (as a numeric depth score of one to four), the ordered list of all hint requests with timestamps, total time in seconds spent at the coding stage, a self-assessed confidence rating from one to five, and the annotations the participant provided at the submission gate. The annotation gate requires the participant to state in free text what their code does and to list which threats they believe they have addressed, before the Code Reviewer is permitted to run.

\subsubsection*{Code Reviewer}

The Code Reviewer receives the participant's code and the threat model from state and conducts an independent security review in two sequential stages. First, it calls the \textit{run\_bandit} MCP tool, which executes Bandit static analysis on the code and returns structured JSON findings. Second, it calls GPT-4o with the code, the formatted Bandit findings, and the threat model mitigations as combined context, instructing the model to include all Bandit findings in the output (tagged \texttt{"source": "bandit"}), add any further issues Bandit missed (tagged \texttt{"source": "llm"}), and verify whether each threat model mitigation was correctly implemented. Each finding in the output is a \textit{ReviewFinding} dictionary with fields \texttt{cwe\_id}, \texttt{severity}, \texttt{description} (referencing specific lines), \texttt{suggested\_fix} (a concrete, implementable fix), \texttt{line\_number}, and \texttt{source}. The raw Bandit findings are stored in \texttt{state["bandit\_findings\_review"]}, which is a separate state field from the Verifier's Bandit run. These two fields are never merged; the Verifier is a genuine independent check, not a repeat.

\subsubsection*{Verifier}

The Verifier performs the final independent validation of the participant's code against the original threat model. It executes four operations. First, it calls the \textit{run\_bandit} MCP tool independently, storing results in \texttt{state["bandit\_findings\_verify"]}. Second, it calls the \textit{execute\_code} MCP tool, which runs the code in a sandboxed subprocess and returns a pass or fail result with exit code, stdout, and stderr. Third, it calls \textit{retrieve(task)} on the CWE corpus independently, with no access to the Threat Modeller's earlier retrieval results, arriving at its own security context. Fourth, it calls GPT-4o with the code, the threat model mitigations, the independent Bandit findings, the independent CWE context, and the sandbox execution result, instructing the model to produce a per-threat verdict (\textit{passed}: true only if the mitigation is fully and correctly implemented, \textit{notes}: a concrete observation citing the specific line or pattern observed). The overall pass verdict is true only when all per-threat verdicts pass and the sandbox execution returns exit code zero.

\subsection{RAG Pipeline}

The Retrieval-Augmented Generation pipeline provides both the Threat Modeller and the Verifier with structured, CWE-specific security context drawn from the MITRE CWE Top 25 Most Dangerous Software Weaknesses corpus. The corpus is stored as a structured JSON file at \textit{rag/cwe\_corpus/cwe\_top25.json}.

Ingestion is handled by \textit{rag/ingest.py}. Each of the 25 CWE entries is split into three focused chunks by \textit{chunk\_entry()}: a main description chunk (containing the CWE identifier, rank, domain, severity, full description, and extended description); a mitigations chunk (containing the prevention strategies for this weakness); and an examples chunk (containing vulnerable code scenarios). This hierarchical splitting ensures that a query about how to prevent a weakness retrieves the mitigations chunk specifically, rather than a general entry in which mitigations constitute only a fraction of the text. The 75 chunks in total are embedded using the OpenAI \texttt{text-embedding-3-small} model and stored in a local ChromaDB persistent collection configured with cosine distance. Each chunk carries metadata fields for \texttt{cwe\_id}, \texttt{rank}, \texttt{severity}, \texttt{domain}, and \texttt{chunk\_type}, which the retriever uses for filtering.

Retrieval is handled by \textit{rag/retriever.py}, which implements the four-stage pipeline. In stage one, \textit{rewrite\_query(task)} calls GPT-4o to expand the coding task into security-relevant search terminology and to identify which of seven available domain labels (\textit{injection}, \textit{memory}, \textit{access\_control}, \textit{input\_validation}, \textit{information\_disclosure}, \textit{network}, \textit{resource\_management}) are relevant to the task. In stage two, ChromaDB's \texttt{where} filter restricts the vector search to chunks whose \texttt{domain} field matches one of the identified domains, reducing the candidate set before any embedding comparison is performed. In stage three, ChromaDB returns the ten closest chunks by cosine similarity to the expanded search query. In stage four, \textit{rerank(task, candidates)} calls GPT-4o to select the three most relevant candidates from the ten returned by the similarity search, reasoning about task fit rather than vector distance alone. The final output is deduplicated by CWE identifier so that the three returned chunks cover distinct weaknesses. The formatted context string returned by \textit{format\_for\_prompt()} is injected directly into each consuming agent's user message.

\subsection{MCP Servers}

The three Model Context Protocol servers are implemented using the \texttt{mcp.server.Server} class from the \texttt{mcp} library, each running over a \texttt{stdio\_server()} transport. They are connected to the agent layer via \textit{mcp\_client.py}, which configures a \textit{MultiServerMCPClient} from \texttt{langchain-mcp-adapters} to manage all three servers as subprocess connections. Each server is launched using the current Python interpreter (\texttt{sys.executable}) so it inherits the project virtual environment. A single \textit{call\_tool(client, tool\_name, arguments)} helper function provides a uniform interface for all agent tool calls, resolving the tool by name across the combined flat tool list and parsing the JSON response into a Python dictionary.

\textit{mcp\_servers/bandit\_server.py} exposes a single tool, \textit{run\_bandit(code)}. The implementation writes the code string to a temporary file using \texttt{tempfile.NamedTemporaryFile}, executes \texttt{bandit -f json -q} on that file as a subprocess with a 30-second timeout, parses the JSON output, normalises each finding into a dictionary with \texttt{test\_id}, \texttt{test\_name}, \texttt{severity}, \texttt{confidence}, \texttt{description}, \texttt{line\_number}, \texttt{cwe\_id}, and \texttt{code\_snippet} fields, and deletes the temporary file in a \texttt{finally} block regardless of whether an error occurred. Bandit exit code zero indicates no issues found; exit code one indicates issues found normally; exit code two indicates that Bandit itself failed to run and is treated as an error.

\textit{mcp\_servers/nist\_nvd\_server.py} exposes a single tool, \textit{search\_nvd(keyword, max\_results)}. The implementation uses \texttt{httpx.AsyncClient} with a ten-second request timeout to query the NVD REST API v2 endpoint at \texttt{services.nvd.nist.gov/rest/json/cves/2.0}, normalises the nested NVD response into a flat list of CVE dictionaries with English description, CVSS v3 severity and score, publication date, associated CWE identifiers, and NVD URL, and returns up to ten results. HTTP 403 and 429 rate-limit responses are handled gracefully with informative error messages. An optional \texttt{NVD\_API\_KEY} environment variable is supported to raise the rolling rate limit from five to fifty requests per thirty seconds.

\textit{mcp\_servers/sandbox\_server.py} exposes a single tool, \textit{execute\_code(code)}. Before spawning any subprocess, the implementation validates the code string with \texttt{ast.parse}, returning a clean syntax-error response immediately if parsing fails. Valid code is written to an isolated temporary directory using \texttt{tempfile.mkdtemp}, which confines the execution working directory to a path outside the project folder so relative file path references cannot reach project files. The code is executed using \texttt{asyncio.create\_subprocess\_exec} with \texttt{stdin=DEVNULL}, integrating the subprocess into the MCP server's event loop and preventing the child process from inheriting the parent's MCP stdio pipe. Execution is bounded by a ten-second timeout via \texttt{asyncio.wait\_for}; if the timeout expires, the process is killed and the \texttt{timed\_out} field in the response is set to true. Stdout and stderr are each capped at 4,096 bytes to prevent memory exhaustion from runaway output. The temporary directory is removed in a \texttt{finally} block.

\subsection{Frontend and Session Management}

The participant-facing interface is a Next.js 15 application using React 19 and TypeScript. Both study views (\textit{frontend/app/study/baseline/page.tsx} and \textit{frontend/app/study/multiagent/page.tsx}) use the Monaco Editor component for code display and editing, loaded asynchronously to avoid blocking the initial page render.

The baseline view manages a five-phase flow: loading (fetching a shuffled task list from the API and reading the participant identifier from session storage), ready (displaying the task and an editor where the participant reads the AI-generated code and may annotate it), generating (awaiting the API response), reviewing (displaying the result with an annotation text area for the participant's notes), and complete (logging the session and advancing to the next task). Task order is randomised server-side on first fetch and the shuffled sequence is stored in session storage so it persists across page refreshes within a session.

The multi-agent view manages a more complex stage-detection flow derived directly from the pipeline state object returned by the API. A \textit{detectStage(state)} function maps the combination of state fields present (plan approved or not, threats approved or not, current\_stage value, review findings present or not) to one of the named UI stages: \textit{plan\_review}, \textit{threat\_review}, \textit{coding}, \textit{code\_review}, \textit{task\_complete}, or \textit{agent\_error}. During the coding stage, the frontend exposes the step plan in a sidebar, calls \textit{fetchStepCaps} to retrieve the per-step hint level caps from the step classifier, and presents hint request controls scoped to each plan step. The \textit{requestStepHint}, \textit{requestNextHint}, and \textit{requestSecurityHint} API calls each invoke the corresponding Code Generator functions. When the participant is ready to submit, the frontend collects their annotation gate responses, confidence rating, and full hint history and calls \textit{finalizeCode}, which resumes the graph and triggers the Code Reviewer. After the Code Reviewer output is displayed, the participant may call \textit{reviseCode} to update their submission before approving the review and releasing the Verifier.

Participant identifiers are assigned at session start, read from URL parameters set by the study coordinator, and stored in browser session storage for the duration of the session. All JSONL log records are written server-side; no sensitive session data is stored in the browser beyond the identifier and task sequence needed to reconstruct the current session state.

%% ============================================================

\section{System Implementation}

\subsection{Overview}

The experiment requires two functionally complete code generation systems. The first, the baseline, is the control condition: a minimal single-agent pipeline in which a participant submits a coding task and receives GPT-4o's response directly, without any intermediate processing. The second, the multi-agent system, is the experimental condition: a five-agent pipeline orchestrated through LangGraph in which each agent performs a discrete, specialised function before the next begins. Both systems use GPT-4o at a fixed temperature of 0.2, are served by the same FastAPI backend, and present to participants through the same Next.js frontend application. The sole variable between conditions is architectural structure. Figure~\ref{fig:architecture} shows the complete deployed system architecture.

\subsection{Baseline Agent}

The baseline agent is implemented in \textit{baseline/agent.py} using the OpenAI Python SDK directly, without any orchestration framework. When a participant submits a task, the \textit{run\_baseline(task, client)} function constructs a two-message prompt containing a brief system message and the task as the user turn, and submits a single completion request to GPT-4o with \texttt{max\_tokens=2048}. The system prompt is deliberately minimal:

\begin{quote}
\textit{You are a software engineer assistant. When given a coding task, produce working, readable code. Do not add commentary beyond what is needed to understand the code. Do not omit error handling.}
\end{quote}

The prompt contains no security instructions, no structured review steps, and no defensive coding guidance. This is intentional: any security guidance in the baseline prompt would reduce the architectural contrast between conditions and bias the comparison. The model response is returned to the participant without further processing.

Session data is written to \textit{logs/baseline\_sessions.jsonl} by \textit{log\_session()}, which appends one JSON object per session containing the fields: \texttt{timestamp} (ISO 8601 UTC), \texttt{condition} (fixed value \texttt{"baseline"}), \texttt{participant\_id}, \texttt{model}, \texttt{temperature}, \texttt{task}, \texttt{task\_id}, \texttt{task\_order}, \texttt{response}, \texttt{duration\_seconds}, and an optional \texttt{participant\_notes} field. The \textit{run\_baseline} function and \textit{log\_session} function are importable module-level callables. The interactive command-line interface for development testing is confined to the \texttt{if \_\_name\_\_ == "\_\_main\_\_"} block and is never imported by the application.

The OpenAI client, model identifier, and temperature are imported from \textit{config.py}, which reads the API key from the environment via \texttt{python-dotenv}. No API credentials are instantiated inside agent files or committed to version control.

\subsection{Multi-Agent System}

The multi-agent system is built on LangGraph's \textit{StateGraph} class, which represents the pipeline as a directed acyclic graph in which each node is an asynchronous function that reads from a shared state object and returns a partial state update. The shared state schema is defined in \textit{multiagent/state.py} as the \textit{AgentState} \textit{TypedDict}, which contains all fields required by all five agents, organised by pipeline stage. Every agent output field is typed as \textit{Optional} because the state is populated progressively and earlier agents cannot know what later agents will produce.

The graph is assembled in \textit{multiagent/graph.py}. Six nodes are registered: \textit{planner}, \textit{threat\_modeller}, \textit{code\_generator}, \textit{code\_reviewer}, \textit{verifier}, and \textit{finalise}. Edges connect them in a linear sequence. The graph is compiled with a \textit{MemorySaver} checkpointer, which persists state across asynchronous resume calls, and with \texttt{interrupt\_after=["planner", "threat\_modeller", "code\_generator", "code\_reviewer"]}. This configuration pauses execution after each of the first four nodes, allowing the frontend to display that node's output to the participant and collect their decision before the pipeline continues.

Three public functions expose the graph to the API layer. \textit{start\_session()} initialises an \textit{AgentState}, generates a unique thread identifier from the participant ID and a UTC timestamp, runs the graph from the entry point through the first interrupt, and returns the thread configuration. \textit{resume\_session()} applies a participant decision as a state update and resumes execution from the next node. \textit{get\_current\_state()} retrieves the latest state snapshot for a given thread, which the frontend polls to read agent output between interrupts.

When the Verifier completes, the graph runs the \textit{\_finalise} node automatically without an additional interrupt. This node computes the total session duration, calls \textit{to\_log\_entry()} to assemble the JSONL record, and appends it to \textit{logs/multiagent\_sessions.jsonl}. The top-level fields of the multi-agent log record (\texttt{timestamp}, \texttt{condition}, \texttt{participant\_id}, \texttt{model}, \texttt{temperature}, \texttt{task}, \texttt{task\_id}, \texttt{task\_order}, \texttt{response}, \texttt{duration\_seconds}) are identical to the baseline schema. An additional \texttt{multi\_agent\_detail} field contains the per-agent outputs, the participant's decision at each interrupt point, and the full HITL coding metrics collected during the code-writing stage.

\subsection{Agent Implementations}

\subsubsection*{Planner}

The Planner receives the participant's coding task from state and produces a structured implementation plan. Its system prompt instructs GPT-4o to decompose the task into three to six concrete, ordered implementation steps, a single-sentence scope statement describing what the code should and should not do, and a list of three to five security properties the final code must satisfy. The prompt explicitly distinguishes concrete security requirements (such as "Hash passwords with bcrypt; do not compare plaintext") from generic advice. The API call uses \texttt{response\_format=\{"type": "json\_object"\}} to guarantee that the response is valid JSON, eliminating the need for error-prone string parsing. The parsed output is stored in \texttt{state["plan"]} as a \textit{PlannerOutput} dictionary. The Planner uses no external tools.

\subsubsection*{Threat Modeller}

The Threat Modeller receives the task and the Planner output and produces a structured threat model before any code is written. It executes three operations in sequence. First, it calls \textit{retrieve(task)} from \textit{rag/retriever.py}, which runs the full four-stage Advanced RAG pipeline described in Section 4.5, returning the three most relevant CWE entries from the local corpus. Second, it queries the NIST NVD REST API v2 via the \textit{search\_nvd} MCP tool, using the short name of the top-ranked CWE entry as the search keyword and requesting three CVEs; the call is wrapped in \texttt{asyncio.wait\_for} with a 25-second timeout, so a slow or unavailable NVD service causes graceful degradation to CWE-only context rather than a pipeline failure. Third, it assembles the task, plan scope, Planner security requirements, formatted CWE context, and formatted CVE context into a single user message and calls GPT-4o to produce a JSON list of threats. Each threat is a \textit{ThreatEntry} dictionary with fields \texttt{cwe\_id}, \texttt{name}, \texttt{severity}, \texttt{description} (task-specific), and \texttt{mitigation} (a direct, implementable instruction for the Code Generator). The raw RAG retrieval results are stored in \texttt{state["rag\_context"]} for logging transparency.

\subsubsection*{Code Generator}

The Code Generator node implements the most distinctive architectural feature of the multi-agent system: a human-in-the-loop coding mentor. In this design, the participant writes all code themselves; no AI-generated code is produced or presented for acceptance. The LangGraph node function \textit{run\_code\_generator()} sets \texttt{state["current\_stage"]} to \texttt{"coding\_in\_progress"} and returns immediately, signalling to the graph that the pipeline is waiting for the human participant rather than completing an autonomous step.

During the coding stage, the frontend provides the participant with four asynchronous hint functions. \textit{get\_step\_hint(state, step\_index, level)} returns a hint for a specific plan step at one of four escalating levels: \textit{direction} (plain English explaining what the step requires and why); \textit{pseudocode} (comment-only skeleton with no executable code); \textit{partial} (structural implementation with \texttt{\# TODO} markers for security-critical logic); and \textit{full} (complete implementation for boilerplate steps). Each level is defined by its own system prompt with strict instructions about what may and may not be revealed. \textit{get\_next\_hint(state, code\_so\_far)} reads the participant's current code and returns a focused suggestion identifying which plan step appears to need attention next, without rewriting any code. \textit{get\_security\_hint(state, code\_so\_far)} scans the current code against the threat model and flags the single most important security issue or missing mitigation.

The maximum hint level available to the participant for each plan step is not uniform. Prior to the coding stage, a step classifier (\textit{multiagent/step\_classifier.py}) uses a separate GPT-4o call at temperature 0.0 to read the plan steps and the threat model mitigations together, assigning each step a cap of \textit{direction}, \textit{pseudocode}, \textit{partial}, or \textit{full} based on whether the step directly implements a named security mitigation, is security-adjacent, is general logic, or is pure boilerplate respectively. Steps classified as \textit{direction} receive only plain-English guidance regardless of what the participant requests; the corresponding code is never revealed. The per-step caps are stored in \texttt{state["step\_hint\_caps"]} and enforced by the API route handler before any hint function is called.

When the participant is ready to submit their code, the frontend calls \textit{finalize\_code(state, user\_code, annotations, confidence\_rating, hints\_requested, time\_in\_coding\_seconds)}, which stores the participant's code in \texttt{state["generated\_code"]} and packages the full HITL metrics into a \textit{HitlCodingMetrics} dictionary: the deepest hint level reached across all steps (as a numeric depth score of one to four), the ordered list of all hint requests with timestamps, total time in seconds spent at the coding stage, a self-assessed confidence rating from one to five, and the annotations the participant provided at the submission gate. The annotation gate requires the participant to state in free text what their code does and to list which threats they believe they have addressed, before the Code Reviewer is permitted to run.

\subsubsection*{Code Reviewer}

The Code Reviewer receives the participant's code and the threat model from state and conducts an independent security review in two sequential stages. First, it calls the \textit{run\_bandit} MCP tool, which executes Bandit static analysis on the code and returns structured JSON findings. Second, it calls GPT-4o with the code, the formatted Bandit findings, and the threat model mitigations as combined context, instructing the model to include all Bandit findings in the output (tagged \texttt{"source": "bandit"}), add any further issues Bandit missed (tagged \texttt{"source": "llm"}), and verify whether each threat model mitigation was correctly implemented. Each finding in the output is a \textit{ReviewFinding} dictionary with fields \texttt{cwe\_id}, \texttt{severity}, \texttt{description} (referencing specific lines), \texttt{suggested\_fix} (a concrete, implementable fix), \texttt{line\_number}, and \texttt{source}. The raw Bandit findings are stored in \texttt{state["bandit\_findings\_review"]}, which is a separate state field from the Verifier's Bandit run. These two fields are never merged; the Verifier is a genuine independent check, not a repeat.

\subsubsection*{Verifier}

The Verifier performs the final independent validation of the participant's code against the original threat model. It executes four operations. First, it calls the \textit{run\_bandit} MCP tool independently, storing results in \texttt{state["bandit\_findings\_verify"]}. Second, it calls the \textit{execute\_code} MCP tool, which runs the code in a sandboxed subprocess and returns a pass or fail result with exit code, stdout, and stderr. Third, it calls \textit{retrieve(task)} on the CWE corpus independently, with no access to the Threat Modeller's earlier retrieval results, arriving at its own security context. Fourth, it calls GPT-4o with the code, the threat model mitigations, the independent Bandit findings, the independent CWE context, and the sandbox execution result, instructing the model to produce a per-threat verdict (\textit{passed}: true only if the mitigation is fully and correctly implemented, \textit{notes}: a concrete observation citing the specific line or pattern observed). The overall pass verdict is true only when all per-threat verdicts pass and the sandbox execution returns exit code zero.

\subsection{RAG Pipeline}

The Retrieval-Augmented Generation pipeline provides both the Threat Modeller and the Verifier with structured, CWE-specific security context drawn from the MITRE CWE Top 25 Most Dangerous Software Weaknesses corpus. The corpus is stored as a structured JSON file at \textit{rag/cwe\_corpus/cwe\_top25.json}.

Ingestion is handled by \textit{rag/ingest.py}. Each of the 25 CWE entries is split into three focused chunks by \textit{chunk\_entry()}: a main description chunk (containing the CWE identifier, rank, domain, severity, full description, and extended description); a mitigations chunk (containing the prevention strategies for this weakness); and an examples chunk (containing vulnerable code scenarios). This hierarchical splitting ensures that a query about how to prevent a weakness retrieves the mitigations chunk specifically, rather than a general entry in which mitigations constitute only a fraction of the text. The 75 chunks in total are embedded using the OpenAI \texttt{text-embedding-3-small} model and stored in a local ChromaDB persistent collection configured with cosine distance. Each chunk carries metadata fields for \texttt{cwe\_id}, \texttt{rank}, \texttt{severity}, \texttt{domain}, and \texttt{chunk\_type}, which the retriever uses for filtering.

Retrieval is handled by \textit{rag/retriever.py}, which implements the four-stage pipeline. In stage one, \textit{rewrite\_query(task)} calls GPT-4o to expand the coding task into security-relevant search terminology and to identify which of seven available domain labels (\textit{injection}, \textit{memory}, \textit{access\_control}, \textit{input\_validation}, \textit{information\_disclosure}, \textit{network}, \textit{resource\_management}) are relevant to the task. In stage two, ChromaDB's \texttt{where} filter restricts the vector search to chunks whose \texttt{domain} field matches one of the identified domains, reducing the candidate set before any embedding comparison is performed. In stage three, ChromaDB returns the ten closest chunks by cosine similarity to the expanded search query. In stage four, \textit{rerank(task, candidates)} calls GPT-4o to select the three most relevant candidates from the ten returned by the similarity search, reasoning about task fit rather than vector distance alone. The final output is deduplicated by CWE identifier so that the three returned chunks cover distinct weaknesses. The formatted context string returned by \textit{format\_for\_prompt()} is injected directly into each consuming agent's user message.

\subsection{MCP Servers}

The three Model Context Protocol servers are implemented using the \texttt{mcp.server.Server} class from the \texttt{mcp} library, each running over a \texttt{stdio\_server()} transport. They are connected to the agent layer via \textit{mcp\_client.py}, which configures a \textit{MultiServerMCPClient} from \texttt{langchain-mcp-adapters} to manage all three servers as subprocess connections. Each server is launched using the current Python interpreter (\texttt{sys.executable}) so it inherits the project virtual environment. A single \textit{call\_tool(client, tool\_name, arguments)} helper function provides a uniform interface for all agent tool calls, resolving the tool by name across the combined flat tool list and parsing the JSON response into a Python dictionary.

\textit{mcp\_servers/bandit\_server.py} exposes a single tool, \textit{run\_bandit(code)}. The implementation writes the code string to a temporary file using \texttt{tempfile.NamedTemporaryFile}, executes \texttt{bandit -f json -q} on that file as a subprocess with a 30-second timeout, parses the JSON output, normalises each finding into a dictionary with \texttt{test\_id}, \texttt{test\_name}, \texttt{severity}, \texttt{confidence}, \texttt{description}, \texttt{line\_number}, \texttt{cwe\_id}, and \texttt{code\_snippet} fields, and deletes the temporary file in a \texttt{finally} block regardless of whether an error occurred. Bandit exit code zero indicates no issues found; exit code one indicates issues found normally; exit code two indicates that Bandit itself failed to run and is treated as an error.

\textit{mcp\_servers/nist\_nvd\_server.py} exposes a single tool, \textit{search\_nvd(keyword, max\_results)}. The implementation uses \texttt{httpx.AsyncClient} with a ten-second request timeout to query the NVD REST API v2 endpoint at \texttt{services.nvd.nist.gov/rest/json/cves/2.0}, normalises the nested NVD response into a flat list of CVE dictionaries with English description, CVSS v3 severity and score, publication date, associated CWE identifiers, and NVD URL, and returns up to ten results. HTTP 403 and 429 rate-limit responses are handled gracefully with informative error messages. An optional \texttt{NVD\_API\_KEY} environment variable is supported to raise the rolling rate limit from five to fifty requests per thirty seconds.

\textit{mcp\_servers/sandbox\_server.py} exposes a single tool, \textit{execute\_code(code)}. Before spawning any subprocess, the implementation validates the code string with \texttt{ast.parse}, returning a clean syntax-error response immediately if parsing fails. Valid code is written to an isolated temporary directory using \texttt{tempfile.mkdtemp}, which confines the execution working directory to a path outside the project folder so relative file path references cannot reach project files. The code is executed using \texttt{asyncio.create\_subprocess\_exec} with \texttt{stdin=DEVNULL}, integrating the subprocess into the MCP server's event loop and preventing the child process from inheriting the parent's MCP stdio pipe. Execution is bounded by a ten-second timeout via \texttt{asyncio.wait\_for}; if the timeout expires, the process is killed and the \texttt{timed\_out} field in the response is set to true. Stdout and stderr are each capped at 4,096 bytes to prevent memory exhaustion from runaway output. The temporary directory is removed in a \texttt{finally} block.

\subsection{Frontend and Session Management}

The participant-facing interface is a Next.js 15 application using React 19 and TypeScript. Both study views (\textit{frontend/app/study/baseline/page.tsx} and \textit{frontend/app/study/multiagent/page.tsx}) use the Monaco Editor component for code display and editing, loaded asynchronously to avoid blocking the initial page render.

The baseline view manages a five-phase flow: loading (fetching a shuffled task list from the API and reading the participant identifier from session storage), ready (displaying the task and an editor where the participant reads the AI-generated code and may annotate it), generating (awaiting the API response), reviewing (displaying the result with an annotation text area for the participant's notes), and complete (logging the session and advancing to the next task). Task order is randomised server-side on first fetch and the shuffled sequence is stored in session storage so it persists across page refreshes within a session.

The multi-agent view manages a more complex stage-detection flow derived directly from the pipeline state object returned by the API. A \textit{detectStage(state)} function maps the combination of state fields present (plan approved or not, threats approved or not, current\_stage value, review findings present or not) to one of the named UI stages: \textit{plan\_review}, \textit{threat\_review}, \textit{coding}, \textit{code\_review}, \textit{task\_complete}, or \textit{agent\_error}. During the coding stage, the frontend exposes the step plan in a sidebar, calls \textit{fetchStepCaps} to retrieve the per-step hint level caps from the step classifier, and presents hint request controls scoped to each plan step. The \textit{requestStepHint}, \textit{requestNextHint}, and \textit{requestSecurityHint} API calls each invoke the corresponding Code Generator functions. When the participant is ready to submit, the frontend collects their annotation gate responses, confidence rating, and full hint history and calls \textit{finalizeCode}, which resumes the graph and triggers the Code Reviewer. After the Code Reviewer output is displayed, the participant may call \textit{reviseCode} to update their submission before approving the review and releasing the Verifier.

Participant identifiers are assigned at session start, read from URL parameters set by the study coordinator, and stored in browser session storage for the duration of the session. All JSONL log records are written server-side; no sensitive session data is stored in the browser beyond the identifier and task sequence needed to reconstruct the current session state.

%% ============================================================

\section{System Implementation}

\subsection{Overview}

The experiment requires two functionally complete code generation systems. The first, the baseline, is the control condition: a minimal single-agent pipeline in which a participant submits a coding task and receives GPT-4o's response directly, without any intermediate processing. The second, the multi-agent system, is the experimental condition: a five-agent pipeline orchestrated through LangGraph in which each agent performs a discrete, specialised function before the next begins. Both systems use GPT-4o at a fixed temperature of 0.2, are served by the same FastAPI backend, and present to participants through the same Next.js frontend application. The sole variable between conditions is architectural structure. Figure~\ref{fig:architecture} shows the complete deployed system architecture.

\subsection{Baseline Agent}

The baseline agent is implemented in \textit{baseline/agent.py} using the OpenAI Python SDK directly, without any orchestration framework. When a participant submits a task, the \textit{run\_baseline(task, client)} function constructs a two-message prompt containing a brief system message and the task as the user turn, and submits a single completion request to GPT-4o with \texttt{max\_tokens=2048}. The system prompt is deliberately minimal:

\begin{quote}
\textit{You are a software engineer assistant. When given a coding task, produce working, readable code. Do not add commentary beyond what is needed to understand the code. Do not omit error handling.}
\end{quote}

The prompt contains no security instructions, no structured review steps, and no defensive coding guidance. This is intentional: any security guidance in the baseline prompt would reduce the architectural contrast between conditions and bias the comparison. The model response is returned to the participant without further processing.

Session data is written to \textit{logs/baseline\_sessions.jsonl} by \textit{log\_session()}, which appends one JSON object per session containing the fields: \texttt{timestamp} (ISO 8601 UTC), \texttt{condition} (fixed value \texttt{"baseline"}), \texttt{participant\_id}, \texttt{model}, \texttt{temperature}, \texttt{task}, \texttt{task\_id}, \texttt{task\_order}, \texttt{response}, \texttt{duration\_seconds}, and an optional \texttt{participant\_notes} field. The \textit{run\_baseline} function and \textit{log\_session} function are importable module-level callables. The interactive command-line interface for development testing is confined to the \texttt{if \_\_name\_\_ == "\_\_main\_\_"} block and is never imported by the application.

The OpenAI client, model identifier, and temperature are imported from \textit{config.py}, which reads the API key from the environment via \texttt{python-dotenv}. No API credentials are instantiated inside agent files or committed to version control.

\subsection{Multi-Agent System}

The multi-agent system is built on LangGraph's \textit{StateGraph} class, which represents the pipeline as a directed acyclic graph in which each node is an asynchronous function that reads from a shared state object and returns a partial state update. The shared state schema is defined in \textit{multiagent/state.py} as the \textit{AgentState} \textit{TypedDict}, which contains all fields required by all five agents, organised by pipeline stage. Every agent output field is typed as \textit{Optional} because the state is populated progressively and earlier agents cannot know what later agents will produce.

The graph is assembled in \textit{multiagent/graph.py}. Six nodes are registered: \textit{planner}, \textit{threat\_modeller}, \textit{code\_generator}, \textit{code\_reviewer}, \textit{verifier}, and \textit{finalise}. Edges connect them in a linear sequence. The graph is compiled with a \textit{MemorySaver} checkpointer, which persists state across asynchronous resume calls, and with \texttt{interrupt\_after=["planner", "threat\_modeller", "code\_generator", "code\_reviewer"]}. This configuration pauses execution after each of the first four nodes, allowing the frontend to display that node's output to the participant and collect their decision before the pipeline continues.

Three public functions expose the graph to the API layer. \textit{start\_session()} initialises an \textit{AgentState}, generates a unique thread identifier from the participant ID and a UTC timestamp, runs the graph from the entry point through the first interrupt, and returns the thread configuration. \textit{resume\_session()} applies a participant decision as a state update and resumes execution from the next node. \textit{get\_current\_state()} retrieves the latest state snapshot for a given thread, which the frontend polls to read agent output between interrupts.

When the Verifier completes, the graph runs the \textit{\_finalise} node automatically without an additional interrupt. This node computes the total session duration, calls \textit{to\_log\_entry()} to assemble the JSONL record, and appends it to \textit{logs/multiagent\_sessions.jsonl}. The top-level fields of the multi-agent log record (\texttt{timestamp}, \texttt{condition}, \texttt{participant\_id}, \texttt{model}, \texttt{temperature}, \texttt{task}, \texttt{task\_id}, \texttt{task\_order}, \texttt{response}, \texttt{duration\_seconds}) are identical to the baseline schema. An additional \texttt{multi\_agent\_detail} field contains the per-agent outputs, the participant's decision at each interrupt point, and the full HITL coding metrics collected during the code-writing stage.

\subsection{Agent Implementations}

\subsubsection*{Planner}

The Planner receives the participant's coding task from state and produces a structured implementation plan. Its system prompt instructs GPT-4o to decompose the task into three to six concrete, ordered implementation steps, a single-sentence scope statement describing what the code should and should not do, and a list of three to five security properties the final code must satisfy. The prompt explicitly distinguishes concrete security requirements (such as "Hash passwords with bcrypt; do not compare plaintext") from generic advice. The API call uses \texttt{response\_format=\{"type": "json\_object"\}} to guarantee that the response is valid JSON, eliminating the need for error-prone string parsing. The parsed output is stored in \texttt{state["plan"]} as a \textit{PlannerOutput} dictionary. The Planner uses no external tools.

\subsubsection*{Threat Modeller}

The Threat Modeller receives the task and the Planner output and produces a structured threat model before any code is written. It executes three operations in sequence. First, it calls \textit{retrieve(task)} from \textit{rag/retriever.py}, which runs the full four-stage Advanced RAG pipeline described in Section 4.5, returning the three most relevant CWE entries from the local corpus. Second, it queries the NIST NVD REST API v2 via the \textit{search\_nvd} MCP tool, using the short name of the top-ranked CWE entry as the search keyword and requesting three CVEs; the call is wrapped in \texttt{asyncio.wait\_for} with a 25-second timeout, so a slow or unavailable NVD service causes graceful degradation to CWE-only context rather than a pipeline failure. Third, it assembles the task, plan scope, Planner security requirements, formatted CWE context, and formatted CVE context into a single user message and calls GPT-4o to produce a JSON list of threats. Each threat is a \textit{ThreatEntry} dictionary with fields \texttt{cwe\_id}, \texttt{name}, \texttt{severity}, \texttt{description} (task-specific), and \texttt{mitigation} (a direct, implementable instruction for the Code Generator). The raw RAG retrieval results are stored in \texttt{state["rag\_context"]} for logging transparency.

\subsubsection*{Code Generator}

The Code Generator node implements the most distinctive architectural feature of the multi-agent system: a human-in-the-loop coding mentor. In this design, the participant writes all code themselves; no AI-generated code is produced or presented for acceptance. The LangGraph node function \textit{run\_code\_generator()} sets \texttt{state["current\_stage"]} to \texttt{"coding\_in\_progress"} and returns immediately, signalling to the graph that the pipeline is waiting for the human participant rather than completing an autonomous step.

During the coding stage, the frontend provides the participant with four asynchronous hint functions. \textit{get\_step\_hint(state, step\_index, level)} returns a hint for a specific plan step at one of four escalating levels: \textit{direction} (plain English explaining what the step requires and why); \textit{pseudocode} (comment-only skeleton with no executable code); \textit{partial} (structural implementation with \texttt{\# TODO} markers for security-critical logic); and \textit{full} (complete implementation for boilerplate steps). Each level is defined by its own system prompt with strict instructions about what may and may not be revealed. \textit{get\_next\_hint(state, code\_so\_far)} reads the participant's current code and returns a focused suggestion identifying which plan step appears to need attention next, without rewriting any code. \textit{get\_security\_hint(state, code\_so\_far)} scans the current code against the threat model and flags the single most important security issue or missing mitigation.

The maximum hint level available to the participant for each plan step is not uniform. Prior to the coding stage, a step classifier (\textit{multiagent/step\_classifier.py}) uses a separate GPT-4o call at temperature 0.0 to read the plan steps and the threat model mitigations together, assigning each step a cap of \textit{direction}, \textit{pseudocode}, \textit{partial}, or \textit{full} based on whether the step directly implements a named security mitigation, is security-adjacent, is general logic, or is pure boilerplate respectively. Steps classified as \textit{direction} receive only plain-English guidance regardless of what the participant requests; the corresponding code is never revealed. The per-step caps are stored in \texttt{state["step\_hint\_caps"]} and enforced by the API route handler before any hint function is called.

When the participant is ready to submit their code, the frontend calls \textit{finalize\_code(state, user\_code, annotations, confidence\_rating, hints\_requested, time\_in\_coding\_seconds)}, which stores the participant's code in \texttt{state["generated\_code"]} and packages the full HITL metrics into a \textit{HitlCodingMetrics} dictionary: the deepest hint level reached across all steps (as a numeric depth score of one to four), the ordered list of all hint requests with timestamps, total time in seconds spent at the coding stage, a self-assessed confidence rating from one to five, and the annotations the participant provided at the submission gate. The annotation gate requires the participant to state in free text what their code does and to list which threats they believe they have addressed, before the Code Reviewer is permitted to run.

\subsubsection*{Code Reviewer}

The Code Reviewer receives the participant's code and the threat model from state and conducts an independent security review in two sequential stages. First, it calls the \textit{run\_bandit} MCP tool, which executes Bandit static analysis on the code and returns structured JSON findings. Second, it calls GPT-4o with the code, the formatted Bandit findings, and the threat model mitigations as combined context, instructing the model to include all Bandit findings in the output (tagged \texttt{"source": "bandit"}), add any further issues Bandit missed (tagged \texttt{"source": "llm"}), and verify whether each threat model mitigation was correctly implemented. Each finding in the output is a \textit{ReviewFinding} dictionary with fields \texttt{cwe\_id}, \texttt{severity}, \texttt{description} (referencing specific lines), \texttt{suggested\_fix} (a concrete, implementable fix), \texttt{line\_number}, and \texttt{source}. The raw Bandit findings are stored in \texttt{state["bandit\_findings\_review"]}, which is a separate state field from the Verifier's Bandit run. These two fields are never merged; the Verifier is a genuine independent check, not a repeat.

\subsubsection*{Verifier}

The Verifier performs the final independent validation of the participant's code against the original threat model. It executes four operations. First, it calls the \textit{run\_bandit} MCP tool independently, storing results in \texttt{state["bandit\_findings\_verify"]}. Second, it calls the \textit{execute\_code} MCP tool, which runs the code in a sandboxed subprocess and returns a pass or fail result with exit code, stdout, and stderr. Third, it calls \textit{retrieve(task)} on the CWE corpus independently, with no access to the Threat Modeller's earlier retrieval results, arriving at its own security context. Fourth, it calls GPT-4o with the code, the threat model mitigations, the independent Bandit findings, the independent CWE context, and the sandbox execution result, instructing the model to produce a per-threat verdict (\textit{passed}: true only if the mitigation is fully and correctly implemented, \textit{notes}: a concrete observation citing the specific line or pattern observed). The overall pass verdict is true only when all per-threat verdicts pass and the sandbox execution returns exit code zero.

\subsection{RAG Pipeline}

The Retrieval-Augmented Generation pipeline provides both the Threat Modeller and the Verifier with structured, CWE-specific security context drawn from the MITRE CWE Top 25 Most Dangerous Software Weaknesses corpus. The corpus is stored as a structured JSON file at \textit{rag/cwe\_corpus/cwe\_top25.json}.

Ingestion is handled by \textit{rag/ingest.py}. Each of the 25 CWE entries is split into three focused chunks by \textit{chunk\_entry()}: a main description chunk (containing the CWE identifier, rank, domain, severity, full description, and extended description); a mitigations chunk (containing the prevention strategies for this weakness); and an examples chunk (containing vulnerable code scenarios). This hierarchical splitting ensures that a query about how to prevent a weakness retrieves the mitigations chunk specifically, rather than a general entry in which mitigations constitute only a fraction of the text. The 75 chunks in total are embedded using the OpenAI \texttt{text-embedding-3-small} model and stored in a local ChromaDB persistent collection configured with cosine distance. Each chunk carries metadata fields for \texttt{cwe\_id}, \texttt{rank}, \texttt{severity}, \texttt{domain}, and \texttt{chunk\_type}, which the retriever uses for filtering.

Retrieval is handled by \textit{rag/retriever.py}, which implements the four-stage pipeline. In stage one, \textit{rewrite\_query(task)} calls GPT-4o to expand the coding task into security-relevant search terminology and to identify which of seven available domain labels (\textit{injection}, \textit{memory}, \textit{access\_control}, \textit{input\_validation}, \textit{information\_disclosure}, \textit{network}, \textit{resource\_management}) are relevant to the task. In stage two, ChromaDB's \texttt{where} filter restricts the vector search to chunks whose \texttt{domain} field matches one of the identified domains, reducing the candidate set before any embedding comparison is performed. In stage three, ChromaDB returns the ten closest chunks by cosine similarity to the expanded search query. In stage four, \textit{rerank(task, candidates)} calls GPT-4o to select the three most relevant candidates from the ten returned by the similarity search, reasoning about task fit rather than vector distance alone. The final output is deduplicated by CWE identifier so that the three returned chunks cover distinct weaknesses. The formatted context string returned by \textit{format\_for\_prompt()} is injected directly into each consuming agent's user message.

\subsection{MCP Servers}

The three Model Context Protocol servers are implemented using the \texttt{mcp.server.Server} class from the \texttt{mcp} library, each running over a \texttt{stdio\_server()} transport. They are connected to the agent layer via \textit{mcp\_client.py}, which configures a \textit{MultiServerMCPClient} from \texttt{langchain-mcp-adapters} to manage all three servers as subprocess connections. Each server is launched using the current Python interpreter (\texttt{sys.executable}) so it inherits the project virtual environment. A single \textit{call\_tool(client, tool\_name, arguments)} helper function provides a uniform interface for all agent tool calls, resolving the tool by name across the combined flat tool list and parsing the JSON response into a Python dictionary.

\textit{mcp\_servers/bandit\_server.py} exposes a single tool, \textit{run\_bandit(code)}. The implementation writes the code string to a temporary file using \texttt{tempfile.NamedTemporaryFile}, executes \texttt{bandit -f json -q} on that file as a subprocess with a 30-second timeout, parses the JSON output, normalises each finding into a dictionary with \texttt{test\_id}, \texttt{test\_name}, \texttt{severity}, \texttt{confidence}, \texttt{description}, \texttt{line\_number}, \texttt{cwe\_id}, and \texttt{code\_snippet} fields, and deletes the temporary file in a \texttt{finally} block regardless of whether an error occurred. Bandit exit code zero indicates no issues found; exit code one indicates issues found normally; exit code two indicates that Bandit itself failed to run and is treated as an error.

\textit{mcp\_servers/nist\_nvd\_server.py} exposes a single tool, \textit{search\_nvd(keyword, max\_results)}. The implementation uses \texttt{httpx.AsyncClient} with a ten-second request timeout to query the NVD REST API v2 endpoint at \texttt{services.nvd.nist.gov/rest/json/cves/2.0}, normalises the nested NVD response into a flat list of CVE dictionaries with English description, CVSS v3 severity and score, publication date, associated CWE identifiers, and NVD URL, and returns up to ten results. HTTP 403 and 429 rate-limit responses are handled gracefully with informative error messages. An optional \texttt{NVD\_API\_KEY} environment variable is supported to raise the rolling rate limit from five to fifty requests per thirty seconds.

\textit{mcp\_servers/sandbox\_server.py} exposes a single tool, \textit{execute\_code(code)}. Before spawning any subprocess, the implementation validates the code string with \texttt{ast.parse}, returning a clean syntax-error response immediately if parsing fails. Valid code is written to an isolated temporary directory using \texttt{tempfile.mkdtemp}, which confines the execution working directory to a path outside the project folder so relative file path references cannot reach project files. The code is executed using \texttt{asyncio.create\_subprocess\_exec} with \texttt{stdin=DEVNULL}, integrating the subprocess into the MCP server's event loop and preventing the child process from inheriting the parent's MCP stdio pipe. Execution is bounded by a ten-second timeout via \texttt{asyncio.wait\_for}; if the timeout expires, the process is killed and the \texttt{timed\_out} field in the response is set to true. Stdout and stderr are each capped at 4,096 bytes to prevent memory exhaustion from runaway output. The temporary directory is removed in a \texttt{finally} block.

\subsection{Frontend and Session Management}

The participant-facing interface is a Next.js 15 application using React 19 and TypeScript. Both study views (\textit{frontend/app/study/baseline/page.tsx} and \textit{frontend/app/study/multiagent/page.tsx}) use the Monaco Editor component for code display and editing, loaded asynchronously to avoid blocking the initial page render.

The baseline view manages a five-phase flow: loading (fetching a shuffled task list from the API and reading the participant identifier from session storage), ready (displaying the task and an editor where the participant reads the AI-generated code and may annotate it), generating (awaiting the API response), reviewing (displaying the result with an annotation text area for the participant's notes), and complete (logging the session and advancing to the next task). Task order is randomised server-side on first fetch and the shuffled sequence is stored in session storage so it persists across page refreshes within a session.

The multi-agent view manages a more complex stage-detection flow derived directly from the pipeline state object returned by the API. A \textit{detectStage(state)} function maps the combination of state fields present (plan approved or not, threats approved or not, current\_stage value, review findings present or not) to one of the named UI stages: \textit{plan\_review}, \textit{threat\_review}, \textit{coding}, \textit{code\_review}, \textit{task\_complete}, or \textit{agent\_error}. During the coding stage, the frontend exposes the step plan in a sidebar, calls \textit{fetchStepCaps} to retrieve the per-step hint level caps from the step classifier, and presents hint request controls scoped to each plan step. The \textit{requestStepHint}, \textit{requestNextHint}, and \textit{requestSecurityHint} API calls each invoke the corresponding Code Generator functions. When the participant is ready to submit, the frontend collects their annotation gate responses, confidence rating, and full hint history and calls \textit{finalizeCode}, which resumes the graph and triggers the Code Reviewer. After the Code Reviewer output is displayed, the participant may call \textit{reviseCode} to update their submission before approving the review and releasing the Verifier.

Participant identifiers are assigned at session start, read from URL parameters set by the study coordinator, and stored in browser session storage for the duration of the session. All JSONL log records are written server-side; no sensitive session data is stored in the browser beyond the identifier and task sequence needed to reconstruct the current session state.
